\section{Continuum Limit and Scaling Considerations}

  \subsection{Discrete-to-Continuum Correspondence}

    The numerical implementation employs a finite cubic lattice with
    Dirichlet boundary conditions.
    Let $h$ denote the lattice spacing and $L_h$ the corresponding
    graph Laplacian.
    In the limit $h \to 0$ with fixed physical domain size,
    the rescaled operator satisfies
    \begin{equation}
      h^{-2} L_h \longrightarrow -\nabla^2
    \end{equation}
    in the sense of operator convergence on suitable function spaces.

    Accordingly, the discrete resolvent converges to the continuum Green operator,
    \begin{equation}
      (L_h + \varepsilon I)^{-1}
      \;\longrightarrow\;
      (-\nabla^2 + \varepsilon)^{-1},
    \end{equation}
    for fixed $\varepsilon > 0$.
    In three dimensions and for $\varepsilon \to 0$,
    the corresponding Green kernel asymptotically reproduces
    the Newtonian form
    \begin{equation}
      G(r) \sim \frac{1}{4\pi r}.
    \end{equation}

    The discrete $1/r$ behavior observed numerically is therefore
    consistent with the continuum limit of Laplace-type operators.

  \subsection{Scaling of Entropy Variation}

    Under a rescaling of spatial coordinates
    $x \mapsto \lambda x$,
    the Laplacian transforms as
    \begin{equation}
      -\nabla^2 \mapsto \lambda^{-2} (-\nabla^2).
    \end{equation}
    Eigenvalues scale accordingly as $\lambda^{-2}$,
    so that the logarithmic determinant transforms as
    \begin{equation}
      \log \det{}' A \mapsto
      \log \det{}' A - 2 N \log \lambda,
    \end{equation}
    where $N$ denotes the number of positive modes
    below a fixed spectral cutoff.

    Although the determinant itself depends on regularization,
    its first variation under localized perturbations,
    \begin{equation}
      \delta S_\Pi = \frac12 \mathrm{Tr}(A^{-1}\delta A),
    \end{equation}
    is scale-covariant and governed by the Green kernel.
    The $1/r$ decay therefore follows from dimensional analysis
    in three spatial dimensions and is not an artifact of discretization.

  \subsection{Interpretation}

    The continuum scaling properties confirm that
    the Newtonian kernel arises from the elliptic nature
    of the operator rather than from lattice-specific effects.
    The discrete simulations reproduce the expected continuum behavior,
    while making explicit the role of finite spatial support
    in generating a macroscopic amplitude shift.

    Within this framework, gravity appears as a long-range
    resolvent-mediated response of a projective entropy functional.
    The derivation is restricted to the weak-field,
    Newtonian regime and does not rely on additional geometric structure.
